%% BEGIN ### RELEASE

\begin{comment}
    ⚠️ Although this is a .tex file, it's symlinked from a05-comparisons.md
\end{comment}

\clearpage
\section{Comparison: ``The Big 4''}

The following comparisons (\autoref{table:compare_nets_3k}, \autoref{table:compare_nets_20k}, and \autoref{table:comparison_1m_tps}) are intended to be an \emph{apples to apples} comparison between UT variants and ``The Big 4'': Bitcoin, Cardano, Eth2, and Polkadot.
Cardano, as a cutting-edge network, is an exception here: the Cardano/IOHK teams have not been pursuing \emph{Layer 1} scalability solutions and are focusing instead on a \emph{Layer 2} solution: Hydra via EUTXOs.
Ethereum is also an exception of sorts: these comparisons use the original \emph{Eth2} sharding roadmap as an example of a sharded Layer 1, whereas the Ethereum network itself is now pursuing a rollup-centric scaling strategy via their beacon chain.
I mention this particularly because the nature of an \emph{apples to apples} comparison casts Cardano and Ethereum in a light that some might consider to be misleading.
However, these comparisons deliberately do not consider \emph{Layer 2} scalability solutions for one very simple reason: all networks can implement them in some fashion.
It is not a fair (or accurate) comparison \emph{of blockchain architecture} if \emph{Layer 2} scaling solutions are considered for some networks and not for others.
\begin{comment}
Additionally, as mentioned in [Ethereum's *Sharding FAQs*](https://web.archive.org/web/20220606205028/https://eth.wiki/sharding/Sharding-FAQs#how-does-plasma-state-channels-and-other-layer-2-technologies-fit-into-the-trilemma), payment- and state-channels provide a *constant factor* increase in throughput, so those techniques will not improve a network's scaling complexity as it is measured here.
\end{comment}

In addition to those four, a network named \emph{Opt.Shard} (for \emph{Optimal Sharding}) is included in these comparisons.
This is a theoretical network which uses the \emph{best parameters possible among UT configurations}.
No real-world sharded network has come close to this level of performance, and it is incompatible with PoS.

\defineTermTex{Scaling Factor}{
    Also: ``Scale $\times$''. For a given $k$, it is the factor by which TPS increases with an additional nesting level.
    In effect, it allows for comparison of the efficacy of scaling schemes when $k$ is fixed.
    For some designs, the \emph{Scaling Factor} can change between nesting levels
}

<!--
\todoDraftOnly[h]{Last row with infinity: theoretically yes, but is it practically achievable?}
mk: added explanatory comment. I checked the code and realized why I put inf there: we need to specify how big the tiling is but that's not bounded.
-->

%% INSERT ### TABLE: compare_nets_3k

: A comparison of quantitative scaling properties between UT and various networks given $k = 3000$ bytes/s. Transaction size is set to 250 bytes, $D_f = B_f$, and $D_h = B_h$.
Note: `$\infty$' here does not literally mean infinite; however, we cannot calculate concrete values without specifying parameters arbitrarily which are otherwise unbounded.

%% INSERT ### TABLE: compare_nets_20k

: Similar to the previous table, but with $k = 20$ kB/s instead of 3 kB/s.

%% INSERT ### TABLE: comparison_1m_tps

: A comparison of computational requirements (approximated by $k$) for 1 million TPS between UT and various other networks. $\UT{2}$'s equivalent TPS is also provided (given the same parameters of $k$, $B_f$ and $B_h$).

\begin{comment}
%% INS--ERT ### TABLE: comparison_1gbps

: A comparison of various networks' $k$ and TPS given a maximum network-wide throughput of 1 Gb/s ($\sim1.3\times 10^8$ B/s). Also included is the rate at which chains on that network grow in size (which is proportional to $k$).
\end{comment}

\ctable[
    pos = htb,
    caption = Table evaluating various other networks against trilemma criteria.,
    cap = Table evaluating various other networks against trilemma criteria.,
    center,
    label = tab:other-nets-trilemma,
]{lllll}{
    \tnote[*]{Real-world performance of Bitcoin; $k \approx 1700$ B/s; $\text{Tx}_\text{avg} \approx 375$ B.}
    \tnote[$\!$\#,$\|$,\P]{Prediction based on $3000 \le k \le 20000$.}
    %\tnote[et]{Prediction based on $3000 \le k \le 20000$.}
    %\tnote[pt]{Prediction based on $3000 \le k \le 20000$.}
    \tnote[\S]{
        Assuming $3000 \le k \le 20000$ and that child-chains use PoW + PoR with the smallest possible headers.
    }
    \tnote[PoS]{
        \href{\linkZackPoSCriticisms}{Unanswered criticisms of PoS} (\href{\linkZackPoSDefence}{replies})
        mean that we cannot conclude that PoS is $O(n)$ secure.
        Saying PoS is ``Maybe'' $O(n)$ secure (especially when used in isolation) reflects an optimistic perspective --- since there are unanswered criticisms, we should arguably conclude ``No''.
        (There are no such unanswered criticisms of PoW-based consensus.)
    }
    \tnote[$\dagger$]{
        \href{https://web.archive.org/web/20241227202028/https://solana.com/solana-whitepaper.pdf}{PoH} is $O(c)$ secure by design.
        Solana uses PoS on top of PoH, but it's unclear which has precedence (and one must).
        Additionally, since a $\le O(c)$ DoS has brought down Solana before, this is an upper limit on Solana's security.
    }
    \tnote[$\ddagger$]{
        Even though Solana requires $k \approx 10^7$ ($\sim$ 12 MB/s) for 50K TPS, this is consistent with their published ``strategy'' for ``scaling''.
        50K is chosen as the limit because \href{\linkSolanaFiftyKClaim}{the official site claims \emph{at least} 50K as the upper limit} and \href{\linkSolanaFourHundKFail}{400K is known to be beyond Solana's capabilities}.
        (Note: this also earns Solana a decisive ``No'' under the ``Decentralized?'' column.)
    }
}{
        \FL
        Network\hspace{2em} & Decentralized?\hspace{1em} & $O(n)$ Secure?\hspace{1em} & \multicolumn{2}{l}{$O(n)$ Scalable?}
        \ML
        Bitcoin   & Yes           & Yes              & No    & (Max. TPS: $\sim$ 5)\tmark[*] \NN
        Cardano   & Yes           & Maybe\tmark[PoS] & No    & (Max. TPS: $\sim$ 10 - 80)\tmark[$\!$\#] \NN
        Solana    & No\tmark[$\ddagger$]  & No ($O(c)$)\tmark[$\dagger$] & Unlikely & (Max. TPS: $\sim$ 50K)\tmark[$\ddagger$] \NN
        Polkadot  & Maybe         & Maybe\tmark[PoS] & Unlikely & (Max. TPS: $\sim$ 200 - 12K)\tmark[\P] \NN
        Eth2      & Maybe         & Maybe\tmark[PoS] & Unlikely & (Max. TPS: $\sim$ 1K - 62K)\tmark[$\|$] \NN
        Opt.Shard & Yes           & Yes              & Maybe & (Max. TPS: $\sim$ 8K - 350K)\tmark[\S]
        \LL
}

%% END ### RELEASE
